% Folio: Starrer Körper (EP1), auf Deutsch. Build: python engine/kl.py build <this file>
\documentclass[deutsch]{kaelearn}
\subject{Experimentalphysik I}
\folio{01}
\title{Starrer\\Körper}
\shorttitle{Starrer Körper}
\subtitle{$10^{25}$ Atome, sechs Zahlen.\\Wie sich ein Ding bewegt, das sich nicht verformt.}
\sources{Deine EP1-Mitschrift · Standardstoff EP1}
\readtime{2 Min. Überblick · 75 Min. Folio}
\filed{2026-11-09}
\plate{plate.png}{Das Orrery}
\kael{reading}

\newcommand\dd{\mathrm{d}}
\newcommand\ext{_{\mathrm{ext}}}
\newcommand\step[1]{\par\smallskip\noindent{\color{kldim}\pix{#1}}\quad}
\makeatletter
\newenvironment{shortexample}[1]{\refstepcounter{klitem}\begin{panel}[unbreakable]{\kl@L@ex\ \kl@folio.\theklitem}{#1}{}}{\end{panel}}
\makeatother
\tikzset{vec/.style={-{Stealth[length=2.2mm]},line width=.8pt},lab/.style={fill=klpage,inner sep=1.2pt}}

\begin{document}
\makecover

% ======================================================================
\section{Sechs Zahlen}

\opening{E}{in Ziegelstein} besteht aus etwa $10^{25}$ Atomen. Newton für jedes einzelne wären $3\cdot10^{25}$ gekoppelte
Differentialgleichungen. Niemand löst die. Und trotzdem weißt du ungefähr, wie ein geworfener Ziegel fliegt.

Der Grund ist eine einzige Beobachtung: Im Ziegel ändert sich der Abstand zwischen zwei Atomen nicht. Diese Einschränkung,
$|\vec r_i-\vec r_j| = \text{const}$ für alle Paare, lässt von den $3\cdot10^{25}$ Zahlen genau sechs übrig: drei für den Ort eines
Punktes, drei für die Orientierung. Ein Körper, für den das exakt gilt, heißt \emph{starrer Körper} \de{(rigid body)}. Echte Körper
verformen sich ein wenig; wenn diese Verformung für deine Frage keine Rolle spielt, rechnest du starr.

Das ganze Kapitel beantwortet eine Frage: \emph{Welche Gleichungen regieren diese sechs Zahlen?} Die Antwort sind zwei
Vektorgleichungen, eine für das Schieben, eine für das Drehen. Alles andere, Trägheitsmoment, Tensor, Steiner, Rollen, Kreisel, ist
die Kunst, den Drehimpuls durch die Drehgeschwindigkeit auszudrücken. Wenn du das einmal gesehen hast, ist das Kapitel kurz.
\aside{Euler hat das 1765 aufgeschrieben, in der \emph{Theoria motus corporum solidorum}, samt Hauptachsen. Er sah damals auf einem Auge. Es hat gereicht.}

\begin{twominute}
\begin{itemize}
  \item \textbf{Starr} heißt: alle Abstände fest. Dann reichen 6 Zahlen: 3 für den Schwerpunkt \de{(centre of mass)}, 3 für die Drehlage.
  \item \textbf{Jede Bewegung} ist Schieben plus Drehen: $\vec v_i = \vec v_s + \vec\omega\times\vec r_{si}$, mit \emph{einem} $\vec\omega$ für den ganzen Körper.
  \item \textbf{Schieben:} $M\dot{\vec v}_s = \vec F\ext$. Der Schwerpunkt fliegt wie ein Massenpunkt, egal wie der Körper sich dreht.
  \item \textbf{Drehen:} $\dot{\vec L} = \vec D\ext$, um den Schwerpunkt oder um einen raumfesten Punkt.
  \item \textbf{Die Brücke:} $\vec L = I\vec\omega$ mit dem Trägheitstensor $I_{ij} = \int(r^2\delta_{ij}-x_ix_j)\,\dd m$.
        Um eine feste Achse: $L_z = I\omega$, $D_z = I\dot\omega$, $I = \int r_\perp^2\,\dd m$.
  \item \textbf{Energie:} $E = \frac12Mv_s^2 + \frac12\vec\omega\cdot I_s\vec\omega$. \textbf{Steiner:} $I_A = I_s + Md^2$, nur von einer Achse durch $S$ aus.
  \item \textbf{Schwerkraft} greift im Schwerpunkt an: $\vec D = M\vec R_s\times\vec g$.
  \item \textbf{$\vec L\parallel\vec\omega$ nur um Hauptachsen} (Eigenvektoren von $I$). Sonst läuft $\vec L$ mit, und die Lager müssen ein Drehmoment liefern.
  \item \textbf{Rollen:} Der Auflagepunkt ruht, $v_s = \omega R$. \textbf{Kreisel:} $\vec D\perp\vec L$ dreht $\vec L$, statt den Kreisel zu kippen; $\Omega = D/L$.
  \item \textbf{Der Trick fast jeder Aufgabe:} Leg den Bezugspunkt dorthin, wo die unbekannte Kraft angreift. Dann fällt sie heraus.
\end{itemize}
\end{twominute}

\begin{guessfirst}
\begin{enumerate}
  \item Vollkugel, Vollzylinder und dünner Ring werden gleichzeitig auf derselben schiefen Ebene losgelassen und rollen hinunter.
        In welcher Reihenfolge kommen sie an? Hängt das von Masse oder Radius ab?
  \item Ein Fahrrad fährt mit $10\,\mathrm{m/s}$. Wie schnell bewegt sich der oberste Punkt des Reifens gegenüber der Straße? Der unterste?
  \item Ein schnell drehendes Rad hängt mit waagrechter Achse an einer Schnur, die nur an \emph{einem} Achsenende befestigt ist. Du
        lässt das andere Ende los. Was passiert?
\end{enumerate}
\end{guessfirst}

\begin{terms}
\term{starrer Körper}{rigid body}{alle Abstände zwischen Massenelementen fest}
\term{Schwerpunkt}{centre of mass}{$\vec R_s = \frac1M\int\vec r\,\dd m$}
\term{Winkel\-geschwindigkeit}{angular velocity}{$\vec\omega$: Richtung = Drehachse (rechte Hand), Betrag = Drehrate}
\term{Drehmoment}{torque}{$\vec D = \vec r\times\vec F$}
\term{Drehimpuls}{angular momentum}{$\vec L = \int\vec r\times\vec v\,\dd m$}
\term{Trägheits\-moment}{moment of inertia}{$I = \int r_\perp^2\,\dd m$: Drehträgheit um eine Achse}
\term{Trägheits\-tensor}{inertia tensor}{die Matrix $I$ mit $\vec L = I\vec\omega$}
\term{Hauptachse}{principal axis}{Achse mit $\vec L\parallel\vec\omega$; Eigenvektor von $I$}
\term{Satz von Steiner}{parallel-axis theorem}{$I_A = I_s + Md^2$}
\term{Präzession}{precession}{Drehung von $\vec L$ unter einem Drehmoment $\perp\vec L$}
\end{terms}

% ======================================================================
\section{Schieben und Drehen}

\opening{B}{evor wir} irgendetwas bewegen, die Buchhaltung. Ein Körper mit Dichte $\rho(\vec r)$ hat die Masse
$M = \int_V\rho(\vec r)\,\dd^3V$, bei konstanter Dichte $M = \rho V$. Jede Summe über Teilchen $\sum_i m_i(\dots)$ wird zum Integral
$\int(\dots)\,\dd m$ mit dem Massenelement $\dd m = \rho\,\dd^3V$. Mehr steckt in den Integralen dieses Kapitels nicht.

\begin{definition}{Schwerpunkt}{centre of mass}
\[
  \vec R_s = \frac1M\sum_i m_i\vec r_i = \frac1M\int_V\vec r\,\rho(\vec r)\,\dd^3V .
\]
Schreib jeden Ort als $\vec r_i = \vec R_s + \vec r_{si}$ ($\vec r_{si}$: Ort relativ zum Schwerpunkt). Dann gilt
\[
  \sum_i m_i\,\vec r_{si} = 0, \qquad\text{kontinuierlich}\quad \int(\vec r-\vec R_s)\,\dd m = 0 .
\]
\end{definition}

Die zweite Zeile ist nur die Definition, umgestellt: $\sum m_i\vec r_{si} = \sum m_i\vec r_i - M\vec R_s = 0$. Merk sie dir trotzdem
gesondert. Sie ist das Arbeitspferd des Kapitels: Jeder Kreuzterm, der unten verschwindet (bei der Schwerkraft, bei der Energie,
beim Drehimpuls, bei Steiner), verschwindet an genau dieser Gleichung. Der Schwerpunkt ist \emph{der} Punkt, um den die Masse
ausbalanciert ist. Darum liegt er auf jeder Symmetrieebene: Spiegelpaare heben sich in $\sum m_i\vec r_{si}$ auf.

\subsection{Warum sechs}
Drei Punkte des Körpers, die nicht auf einer Geraden liegen, legen ihn fest: Jeder weitere Punkt ist durch seine Abstände zu diesen
dreien bestimmt. Drei Punkte haben $9$ Koordinaten, ihre drei gegenseitigen Abstände sind fest, bleiben $9-3 = 6$ Freiheitsgrade
\de{(degrees of freedom)}. Üblich: drei für $\vec R_s$, drei Winkel für die Drehlage.

\subsection{Die Geschwindigkeit jedes Punktes}
\aside{$v = \omega r_\perp$ ist die Unterschrift der Starrheit. Galaxien haben sie nicht: Ihre äußeren Sterne laufen fast so schnell wie die inneren. Ich habe nachgezählt.}Leite $\vec r_i = \vec R_s + \vec r_{si}$ ab: $\vec v_i = \vec v_s + \dot{\vec r}_{si}$. Starr heißt $|\vec r_{si}|^2 = \text{const}$, also
\[
  \frac{\dd}{\dd t}\,\vec r_{si}^{\,2} = 2\,\vec r_{si}\cdot\dot{\vec r}_{si} = 0 \quad\Longrightarrow\quad \dot{\vec r}_{si}\perp\vec r_{si}.
\]
Jeder Punkt bewegt sich senkrecht zu seiner Verbindung mit $S$: Er läuft auf einer Kugel um $S$. Das ist noch nicht ganz der Satz,
den du in deiner Mitschrift hast. Der braucht ein Stück mehr.

\begin{theorem}{Geschwindigkeiten im starren Körper}{velocity field}
Es gibt einen Vektor $\vec\omega(t)$, die \emph{Winkelgeschwindigkeit}, für den an jedem Punkt des Körpers gilt
\[
  \vec v_i = \vec v_s + \vec\omega\times\vec r_{si}.
\]
Schieben des Schwerpunkts plus Drehen um eine Achse durch ihn. Liegt der Körper nicht auf einer Geraden, ist $\vec\omega$ eindeutig.
\end{theorem}

\begin{proof}[Beweisskizze]
Starr heißt auch $(\vec r_{si}-\vec r_{sj})^2 = \text{const}$ für jedes Paar. Ableiten und $\vec r_{si}\cdot\dot{\vec r}_{si} = 0$
benutzen gibt $\vec r_{si}\cdot\dot{\vec r}_{sj} = -\,\vec r_{sj}\cdot\dot{\vec r}_{si}$. Die Zuordnung $\vec r_{si}\mapsto\dot{\vec r}_{si}$
setzt sich deshalb zu einer linearen Abbildung $A$ mit $\vec a\cdot A\vec b = -\vec b\cdot A\vec a$ fort (dass diese Fortsetzung
wohldefiniert ist, überspringe ich). $A$ ist also eine antisymmetrische $3\times3$-Matrix, und jede solche ist ein Kreuzprodukt:
\[
  A = \begin{pmatrix}0&-\omega_3&\omega_2\\ \omega_3&0&-\omega_1\\ -\omega_2&\omega_1&0\end{pmatrix},\qquad A\vec b = \vec\omega\times\vec b .
\]
\end{proof}

\why{Warum reicht $|\vec r_{si}| = \text{const}$ nicht für ein gemeinsames $\vec\omega$?}{Es sagt nur, dass jeder Punkt auf
\emph{seiner} Kugel um $S$ bleibt. Zwei Punkte dürften dann um verschiedene Achsen kreisen, wie zwei Monde um einen Planeten. Dann
ändert sich aber ihr Abstand zueinander. Erst die festen Abstände \emph{zwischen} den Punkten kleben alle Drehungen zu einer einzigen
zusammen.}

\begin{center}
\begin{tikzpicture}[scale=1.25]
  \def\blob{plot[smooth cycle,tension=.8] coordinates {(2.6,1.5)(3.4,0.9)(5.0,1.1)(6.6,2.3)(6.5,3.8)(5.0,4.1)(3.6,3.5)(2.7,2.6)}}
  \fill[pattern={dither25},pattern color=kllow] \blob;
  \draw[kldim,line width=.6pt] \blob;
  \coordinate (O) at (0,0); \coordinate (S) at (4,2.2); \coordinate (P) at (5.7,2.0);
  \draw[vec,klink] (O)--(S) node[midway,above left,lab]{$\vec R_s$};
  \draw[vec,klink] (O)--(P) node[midway,below right,lab]{$\vec r_i$};
  \draw[vec,klhi] (S)--(P) node[midway,below,lab]{$\vec r_{si}$};
  \fill[klink] (O) circle (1.6pt) node[below left]{$O$};
  \fill[klhi] (S) circle (2pt) node[below left,lab]{$S$};
  \fill[klhi] (P) circle (1.6pt);
  \draw[klhi,line width=.6pt] (3.3,2.9) circle (4.5pt); \fill[klhi] (3.3,2.9) circle (1.3pt);
  \node[left=5pt,lab] at (3.3,2.9) {$\vec\omega$ (zu dir)};
  \draw[vec,kldim] (P)--++(1.2,0.3) node[right]{$\vec v_s$};
  \draw[vec,kldim] (P)--++(0.11,0.935) node[left,lab]{$\vec\omega\times\vec r_{si}$};
  \draw[kldim,dashed,line width=.4pt] ($(P)+(1.2,0.3)$)--++(0.11,0.935) ($(P)+(0.11,0.935)$)--++(1.2,0.3);
  \draw[vec,klhi,line width=1.1pt] (P)--++(1.31,1.235) node[above]{$\vec v_i$};
\end{tikzpicture}
\end{center}
\plotcaption{Jeder Punkt bekommt dieselbe Schwerpunktsgeschwindigkeit $\vec v_s$ plus einen Drehanteil, der senkrecht auf $\vec r_{si}$
steht und mit dem Abstand zur Achse wächst. $\vec\omega$ zeigt hier aus dem Blatt: Drehung gegen den Uhrzeigersinn.}

Was $\vec\omega$ bedeutet: Daumen der rechten Hand in Richtung $\vec\omega$, die Finger zeigen die Drehrichtung. Der Betrag
$|\vec\omega\times\vec r| = \omega\,r\sin\theta = \omega\,r_\perp$ ist $\omega$ mal Abstand von der Drehachse. Punkte auf der Achse
haben keinen Drehanteil.

\why{Warum hängt $\vec\omega$ nicht davon ab, welchen Punkt du als Zentrum nimmst?}{Nimm statt $S$ irgendeinen Punkt $P$ des Körpers.
Aus $\vec v_P = \vec v_s + \vec\omega\times\vec r_{sP}$ und $\vec v_i = \vec v_s + \vec\omega\times\vec r_{si}$ folgt durch Subtrahieren
$\vec v_i = \vec v_P + \vec\omega\times(\vec r_i-\vec r_P)$. Gleiches $\vec\omega$, anderer Schiebeanteil. Die Drehung gehört dem
Körper; die Translation hängt von deiner Wahl ab. Das nutzt du beim Rollen aus: Wähl den Punkt, der gerade ruht, und die ganze Bewegung
ist eine reine Drehung um ihn.}

% ======================================================================
\section{Wo die Schwerkraft angreift}

\opening{D}{eine Mitschrift} beginnt hier, und zu Recht: Die erste Frage an einen ausgedehnten Körper ist, wo seine Gewichtskraft
eigentlich angreift. Sie greift überall an, an jedem $\dd m$. Summier das auf, im homogenen Feld $\vec g = -g\,\hat e_z$:
\[
  \vec F = \int\vec g\,\dd m = M\vec g,\qquad
  \vec D_O = \int\vec r\times\vec g\,\dd m = \Big(\int\vec r\,\dd m\Big)\times\vec g = M\,\vec R_s\times\vec g .
\]
Der Trick in der zweiten Rechnung: $\vec g$ ist für alle $\dd m$ derselbe Vektor und darf aus dem Integral. Was übrig bleibt,
$\int\vec r\,\dd m$, ist per Definition $M\vec R_s$.

\begin{theorem}{Schwerkraft auf einen starren Körper}{gravity on a rigid body}
Im homogenen Schwerefeld wirken alle Gewichtskräfte zusammen wie \emph{eine} Kraft $M\vec g$, die im Schwerpunkt angreift:
$\vec F = M\vec g$ und $\vec D_O = M\,\vec R_s\times\vec g$ für jeden Bezugspunkt $O$.
\end{theorem}

Daher der Name Schwer-punkt. Die Voraussetzung \emph{homogen} ist echt: Im ungleichmäßigen Feld (der Mond zieht an der erdnahen Seite
stärker) gibt es keinen solchen Punkt, und das resultierende Drehmoment ist der Grund, warum die Erdachse präzediert.

Was daraus folgt, steht in deiner Mitschrift; hier mit Begründung:
\begin{itemize}
  \item Um den Schwerpunkt selbst ($\vec R_s = 0$) ist $\vec D = 0$. Ein im Schwerpunkt aufgehängter Körper ist drehmomentfrei und
        bleibt in jeder Lage stehen: indifferentes Gleichgewicht.
  \item Aufgehängt in $A$: $\vec D_A = M(\vec R_s-\vec r_A)\times\vec g$. Null genau dann, wenn $S$ senkrecht unter oder über $A$ liegt.
        Liegt $S$ unter $A$ (dein „$\vec g\parallel\vec R_s$“, mit $\vec R_s$ von $A$ aus gemessen), dreht jede Auslenkung zurück: stabil.
        Liegt $S$ über $A$, dreht sie weiter weg: labil.
  \item Praktisch: Häng einen Körper nacheinander an zwei Punkten auf. Die beiden Lotlinien schneiden sich im Schwerpunkt.
\end{itemize}

\begin{center}
\begin{tikzpicture}[scale=1.12]
  % stabil
  \begin{scope}
    \begin{scope}[rotate around={18:(0,0)}]
      \fill[pattern={dither25},pattern color=kllow] (-.5,-2.0) rectangle (.5,.35); \draw[kldim] (-.5,-2.0) rectangle (.5,.35);
      \coordinate (S1) at (0,-.95);
    \end{scope}
    \draw[kldim,dashed] (0,.6)--(0,-2.3);
    \draw[klhi,fill=klpage,line width=.8pt] (0,0) circle (3.2pt); \node[right=5pt,lab] at (0,.05){$A$};
    \node[klhi] at (S1){$\times$}; \draw[vec,klhi] (S1)--++(0,-.7) node[right,lab]{$M\vec g$};
    \draw[vec,klink] (-62:2.3) arc (-62:-100:2.3);
    \node[kldim] at (0,-3.2){\pix{stabil}};
  \end{scope}
  % indifferent
  \begin{scope}[xshift=4cm]
    \begin{scope}[rotate around={30:(0,-.8)}]
      \fill[pattern={dither25},pattern color=kllow] (-.5,-2.0) rectangle (.5,.35); \draw[kldim] (-.5,-2.0) rectangle (.5,.35);
    \end{scope}
    \draw[klhi,fill=klpage,line width=.8pt] (0,-.8) circle (3.2pt); \node[right=5pt,lab] at (0,-.8){$A = S$};
    \node[kldim] at (0,-3.2){\pix{indifferent}};
  \end{scope}
  % labil
  \begin{scope}[xshift=8cm,yshift=-1.6cm]
    \begin{scope}[rotate around={-18:(0,0)}]
      \fill[pattern={dither25},pattern color=kllow] (-.5,-.35) rectangle (.5,2.0); \draw[kldim] (-.5,-.35) rectangle (.5,2.0);
      \coordinate (S3) at (0,.95);
    \end{scope}
    \draw[kldim,dashed] (0,-.4)--(0,2.3);
    \draw[klhi,fill=klpage,line width=.8pt] (0,0) circle (3.2pt); \node[left=5pt,lab] at (0,0){$A$};
    \node[klhi] at (S3){$\times$}; \draw[vec,klhi] (S3)--++(0,-.7) node[right,lab]{$M\vec g$};
    \draw[vec,klink] (80:2.3) arc (80:50:2.3);
    \node[kldim] at (0,-1.6){\pix{labil}};
  \end{scope}
\end{tikzpicture}
\end{center}
\plotcaption{Ausgelenkt um ein paar Grad. Links zieht $M\vec g$ den Schwerpunkt ($\times$) zurück unter $A$, rechts weiter weg von der
Senkrechten über $A$. In der Mitte hat $M\vec g$ keinen Hebelarm.}

\subsection{Gleichgewicht}

\begin{theorem}{Statisches Gleichgewicht}{static equilibrium}
Ein ruhender starrer Körper bleibt in Ruhe genau dann, wenn $\sum\vec F\ext = 0$ und $\sum\vec D\ext = 0$.
Ist $\sum\vec F\ext = 0$, dann ist $\sum\vec D\ext$ für jeden Bezugspunkt derselbe.
\end{theorem}

Der zweite Satz in einer Zeile: $\vec D_P = \sum_i(\vec r_i-\vec r_P)\times\vec F_i = \vec D_O - \vec r_P\times\sum_i\vec F_i$.

\why{Warum darfst du den Bezugspunkt frei wählen, und warum ist das schon die halbe Lösung?}{Wenn $\sum\vec F = 0$, hängt $\vec D_P$
nicht von $P$ ab. Also legst du $P$ dorthin, wo die Kräfte angreifen, die du nicht kennst und nicht brauchst: Ihr Hebelarm ist null, sie
fallen aus der Momentengleichung heraus. Aus drei Unbekannten in einer Gleichung wird eine.}

\begin{example}{Leiter an glatter Wand}
Eine Leiter (Länge $\ell$, Masse $m$, Schwerpunkt in der Mitte) lehnt unter dem Winkel $\alpha$ zum Boden an einer reibungsfreien
Wand. Haftreibungszahl am Boden: $\mu$. Ab welchem Winkel rutscht sie?
\begin{center}
\begin{tikzpicture}[scale=.9]
  \fill[pattern={dither25},pattern color=kllow] (-1.2,-.25) rectangle (3.0,0); \draw[kldim] (-1.2,0)--(2.6,0);
  \fill[pattern={dither25},pattern color=kllow] (2.6,-.25) rectangle (2.85,3.9); \draw[kldim] (2.6,0)--(2.6,3.9);
  \draw[klink,line width=1.8pt] (0,0)--(2.6,3.4);
  \draw[kldim] (.6,0) arc (0:52.6:.6); \node[kldim] at (.85,.3){$\alpha$};
  \draw[vec,klhi] (0,0)--(0,1.3) node[above]{$N_1$};
  \draw[vec,klhi] (0,.08)--(.9,.08); \node[klhi] at (1.05,-.25){$f$};
  \draw[vec,klhi] (2.6,3.4)--(1.6,3.4) node[left]{$N_2$};
  \draw[vec,klhi] (1.3,1.7)--(1.3,.65) node[below right=-1pt,lab]{$mg$};
\end{tikzpicture}
\end{center}
\step{Verstehen} Kräfte: Gewicht $mg$ in der Mitte; Boden: Normalkraft $N_1$ nach oben und Haftreibung $f$ zur Wand hin; glatte
Wand: nur $N_2$, waagrecht. Gesucht: das kleinste $\alpha$ mit $f\le\mu N_1$.
\step{Plan} Gleichgewicht: zwei Kraftgleichungen, eine Momentengleichung, drei Unbekannte. Bezugspunkt: der Fuß, denn dort greifen zwei
der drei Unbekannten an.
\step{Rechnung} Senkrecht: $N_1 = mg$. Waagrecht: $f = N_2$. Momente um den Fuß (Kraft mal senkrechter Abstand):
$N_2\,\ell\sin\alpha = mg\,\frac\ell2\cos\alpha$, also $N_2 = \dfrac{mg}{2\tan\alpha}$. Haften verlangt $f\le\mu N_1$:
\[
  \frac{mg}{2\tan\alpha}\le\mu\,mg \quad\Longleftrightarrow\quad \tan\alpha\ge\frac1{2\mu}.
\]
\step{Rückschau} $\mu\to0$ gibt $\alpha\to90^\circ$: Ohne Reibung steht nur die senkrechte Leiter. Großes $\mu$: fast jeder Winkel
hält. Masse und Länge fallen heraus, was sein muss, denn $\tan\alpha$ und $\mu$ sind reine Zahlen.
\end{example}

\begin{shortexample}{Schwerpunkt einer Halbkreisscheibe}
Homogene Halbkreisscheibe, Radius $R$, Flächendichte $\sigma$, gerader Rand auf der $x$-Achse. Wegen der Spiegelsymmetrie ist
$x_s = 0$. Polarkoordinaten: $\dd A = r\,\dd r\,\dd\varphi$, $y = r\sin\varphi$, $M = \sigma\pi R^2/2$:
\[
  y_s = \frac1M\int y\,\sigma\,\dd A = \frac\sigma M\int_0^R\!r^2\,\dd r\int_0^\pi\!\sin\varphi\,\dd\varphi
      = \frac\sigma M\cdot\frac{R^3}3\cdot2 = \frac{4R}{3\pi}\approx0{,}42\,R .
\]
\step{Rückschau} Kleiner als $R/2$: Nahe dem geraden Rand ist die Scheibe am breitesten, dort sitzt mehr Masse. Das $r$ im Flächenelement
ist der Faktor, den man vergisst.
\end{shortexample}

% ======================================================================
\section{Eine feste Achse}

\opening{B}{evor der} Tensor kommt, der Fall, der die meisten Aufgaben abdeckt: Der Körper dreht sich um eine raumfeste Achse. Rad,
Tür, Rolle, Pendel. Leg die Achse auf $z$. Jeder Punkt läuft auf einem Kreis mit Radius $r_\perp$ (Abstand zur Achse) und hat die
Geschwindigkeit $v = \omega r_\perp$. Die kinetische Energie:
\[
  E_{\mathrm{rot}} = \sum_i\tfrac12m_iv_i^2 = \tfrac12\Big(\sum_i m_ir_{\perp i}^2\Big)\,\omega^2 = \tfrac12I\omega^2 .
\]

\begin{definition}{Trägheitsmoment}{moment of inertia}
Das Trägheitsmoment um eine Achse ist $\displaystyle I = \sum_i m_ir_{\perp i}^2 = \int r_\perp^2\,\dd m$, mit $r_\perp$ = Abstand von
der Achse.
\end{definition}

Dieselbe Zahl steckt im Drehimpuls. Mit $\vec v = \omega(-y,x,0)$ ist $(\vec r\times\vec v)_z = xv_y-yv_x = \omega(x^2+y^2) = \omega r_\perp^2$,
also $L_z = I\omega$. Weil sich $r_\perp$ beim Drehen um $z$ nicht ändert, ist $I$ konstant, und $\dot L_z = I\dot\omega$.

\begin{theorem}{Drehung um eine feste Achse}{fixed-axis rotation}
Dreht sich ein starrer Körper mit $\omega$ um eine raumfeste Achse $z$, dann gilt mit dem Trägheitsmoment $I$ um \emph{genau diese} Achse
\[
  L_z = I\omega,\qquad D_z = I\dot\omega,\qquad E_{\mathrm{rot}} = \tfrac12I\omega^2,\qquad P = D_z\omega,\qquad W = \int D_z\,\dd\varphi .
\]
\end{theorem}

Achtung, nur die $z$-Komponente. $L_x$ und $L_y$ müssen nicht null sein; um sie kümmert sich Teil V.

\begin{panel}[unbreakable]{Tabelle}{Schieben und Drehen}{every formula has a twin}
\renewcommand\arraystretch{1.25}
\begin{tabular}{@{}p{34mm}p{36mm}l@{}}
  {\color{kldim}\pix{Größe}} & {\color{kldim}\pix{Translation}} & {\color{kldim}\pix{Drehung um feste Achse}}\\
  Lage & $x$ & $\varphi$\\
  Geschwindigkeit & $v = \dot x$ & $\omega = \dot\varphi$\\
  Beschleunigung & $a = \dot v$ & $\dot\omega$\\
  Trägheit & $m$ & $I = \int r_\perp^2\,\dd m$\\
  Ursache & Kraft $F$ & Drehmoment $D_z = (\vec r\times\vec F)_z$\\
  Impuls & $p = mv$ & $L_z = I\omega$\\
  Grundgleichung & $F = \dot p = ma$ & $D_z = \dot L_z = I\dot\omega$\\
  Energie & $\frac12mv^2$ & $\frac12I\omega^2$\\
  Arbeit, Leistung & $\int F\,\dd x$, \ $Fv$ & $\int D_z\,\dd\varphi$, \ $D_z\omega$
\end{tabular}
\end{panel}

\why{Warum geht der Abstand zur Achse im Quadrat ein?}{Ein Faktor $r_\perp$ kommt aus der Geschwindigkeit: doppelt so weit draußen,
doppelt so schnell. Der zweite kommt, je nachdem wohin du schaust, aus dem Quadrat in $\frac12mv^2$ oder aus dem Hebelarm in
$\vec r\times\vec v$. Probier es: Halt einen Hammer am Stielende und schwing ihn hin und her, dann halt ihn am Kopf. Gleiche Masse,
anderes $I$, ganz anderes Gefühl. Trägheitsmoment misst nicht, \emph{wie viel} Masse da ist, sondern \emph{wo}.}

\subsection{Trägheitsmomente ausrechnen}
Immer dasselbe Rezept: $\dd m$ in passenden Koordinaten ($\lambda\,\dd x$ für einen Stab, $\sigma\,\dd A$ für eine Platte,
$\rho\,\dd^3V$ für Volumen), $r_\perp$ zur gewählten Achse, integrieren, mit $M$ ausdrücken.

\begin{shortexample}{Dünner Stab}
Länge $\ell$, $\lambda = M/\ell$, Achse quer durch die Mitte: $\displaystyle I = \int_{-\ell/2}^{\ell/2}x^2\,\frac M\ell\,\dd x = \frac{M\ell^2}{12}$.
Achse durchs Ende: $\displaystyle\int_0^\ell x^2\,\frac M\ell\,\dd x = \frac{M\ell^2}3$, viermal so viel, weil jetzt Masse doppelt so weit
draußen sitzen kann.
\end{shortexample}

\begin{example}{Vollzylinder, deine Rechnung}
Zylinderkoordinaten $x = r\cos\varphi$, $y = r\sin\varphi$, $z$, mit $\dd^3V = r\,\dd r\,\dd\varphi\,\dd z$. Achse = Symmetrieachse:
\[
  I_{zz} = \rho\int(x^2+y^2)\,\dd^3V = \rho\int_0^h\!\dd z\int_0^{2\pi}\!\dd\varphi\int_0^R\!r^2\cdot r\,\dd r
         = \rho\,h\,2\pi\,\frac{R^4}4 = \tfrac12\underbrace{\rho\,\pi R^2h}_{M}\,R^2 = \tfrac12MR^2 .
\]
Ein $r^2$ ist $r_\perp^2$, das dritte $r$ ist das Volumenelement. Ohne es käme $\frac13$ statt $\frac12$ heraus. Die Höhe $h$ fällt weg:
Eine Münze und eine Walze mit gleichem $M$ und $R$ haben dasselbe $I_{zz}$.
\end{example}

\begin{example}{Kugel, ohne Kugelkoordinaten}
Addiere die Trägheitsmomente um die drei Koordinatenachsen durch den Mittelpunkt:
\[
  I_x+I_y+I_z = \int\big[(y^2+z^2)+(x^2+z^2)+(x^2+y^2)\big]\,\dd m = 2\int r^2\,\dd m .
\]
Bei einer Kugel sind die drei gleich, also $I = \frac23\int r^2\,\dd m$. Vollkugel: in Schalen zerlegen, $\dd m = \rho\,4\pi r^2\,\dd r$,
$\int r^2\,\dd m = 4\pi\rho R^5/5 = \frac35MR^2$ (mit $M = \frac43\pi R^3\rho$), also $I = \frac25MR^2$. Dünne Hohlkugel: alle Masse bei
$r = R$, $\int r^2\,\dd m = MR^2$, also $I = \frac23MR^2$. Symmetrie vor Algebra.
\end{example}

\begin{panel}[unbreakable]{Tabelle}{Trägheitsmomente}{homogene Körper, Masse M}
\renewcommand\arraystretch{1.3}
\begin{tabular}{@{}p{46mm}p{38mm}l@{}}
  {\color{kldim}\pix{Körper}} & {\color{kldim}\pix{Achse}} & {\color{kldim}\pix{Trägheitsmoment}}\\
  Punktmasse & im Abstand $r$ & $Mr^2$\\
  Dünner Ring, dünner Hohlzylinder & Symmetrieachse & $MR^2$\\
  Vollzylinder, Scheibe & Symmetrieachse & $\frac12MR^2$\\
  Dicker Hohlzylinder & Symmetrieachse & $\frac12M(R_1^2+R_2^2)$\\
  Scheibe & Durchmesser & $\frac14MR^2$\\
  Vollzylinder, Höhe $h$ & quer durch $S$ & $M\big(\frac{R^2}4+\frac{h^2}{12}\big)$\\
  Dünner Stab, Länge $\ell$ & quer durch die Mitte & $\frac1{12}M\ell^2$\\
  Dünner Stab, Länge $\ell$ & quer durch ein Ende & $\frac13M\ell^2$\\
  Vollkugel & durch die Mitte & $\frac25MR^2$\\
  Dünne Hohlkugel & durch die Mitte & $\frac23MR^2$\\
  Quader $a\times b\times c$ & parallel zu $c$, durch $S$ & $\frac1{12}M(a^2+b^2)$\\
  Würfel, Kante $a$ & jede Achse durch $S$ & $\frac16Ma^2$
\end{tabular}
\end{panel}

\aside{Huygens hatte den Satz 1673, für seine Pendeluhr. Der deutsche Name ging an Steiner, 150 Jahre später. So ist das mit Namen.}Zwei Abkürzungen erzeugen die halbe Tabelle. Für eine \emph{dünne ebene Platte} in der $xy$-Ebene ist $z = 0$ für alle Masse, also
$I_z = \int(x^2+y^2)\,\dd m = I_x+I_y$ (Satz der senkrechten Achsen). Die Scheibe: $I_z = \frac12MR^2$ und $I_x = I_y$ geben
$I_x = \frac14MR^2$. Die zweite Abkürzung ist Steiner.

\begin{theorem}{Satz von Steiner}{parallel-axis theorem}
Ist $I_s$ das Trägheitsmoment um eine Achse durch den Schwerpunkt und $A$ eine dazu \emph{parallele} Achse im Abstand $d$, dann
\[
  I_A = I_s + Md^2 .
\]
\end{theorem}
\begin{proof}
Miss alles senkrecht zur Achsenrichtung. Der Abstandsvektor eines Massenelements von der Achse $A$ ist $\vec d + \vec u$, wobei $\vec u$ der
Abstandsvektor von der Schwerpunktsachse ist. Dann
$I_A = \int|\vec d+\vec u|^2\,\dd m = Md^2 + 2\,\vec d\cdot\!\int\vec u\,\dd m + \int u^2\,\dd m$.
Der mittlere Term ist null: $\int\vec u\,\dd m$ ist der senkrechte Teil von $\int(\vec r-\vec R_s)\,\dd m = 0$. Das Arbeitspferd, zum
zweiten Mal.
\end{proof}

Probe: Stab ums Ende, $\frac1{12}M\ell^2 + M(\frac\ell2)^2 = \frac13M\ell^2$. Zwei Folgen. Unter allen parallelen Achsen hat die durch $S$
das kleinste $I$. Und Steiner verbindet nur eine Schwerpunktsachse mit einer anderen Achse: Von $A$ nach $B$ gehst du über $S$,
$I_B = I_A - Md_A^2 + Md_B^2$, niemals mit dem Abstand $AB$.

% ======================================================================
\section{Wenn \texorpdfstring{$L$}{L} nicht zu \texorpdfstring{$\omega$}{omega} passt}

\opening{B}{is jetzt} haben wir nur $L_z$ angeschaut. Aber $\vec L$ hat drei Komponenten, und die zwei anderen sind nicht immer null.
Das konkreteste Beispiel: eine Hantel, schräg auf eine Achse gesteckt.

Zwei Massen $m$ sitzen bei $\pm\vec r$, $|\vec r| = a$, der Stab schließt mit der Achse den Winkel $\theta$ ein, und $\vec\omega = \omega\hat e_z$.
In dem Moment, in dem der Stab in der $xz$-Ebene liegt, ist $\vec r = a(\sin\theta,0,\cos\theta)$ und $\vec v = \vec\omega\times\vec r = \omega a\sin\theta\,\hat e_y$.
Beide Massen zusammen:
\[
  \vec L = 2m\,\vec r\times\vec v = 2m\,\omega a^2\sin\theta\,(-\cos\theta,\ 0,\ \sin\theta).
\]
Dieser Vektor steht senkrecht auf dem Stab (Skalarprodukt mit $(\sin\theta,0,\cos\theta)$ ist null), nicht parallel zur Achse. Seine
$z$-Komponente ist $2m(a\sin\theta)^2\omega = I_z\omega$, wie in Teil IV. Aber er hat auch $L_x = -2m\omega a^2\sin\theta\cos\theta$, und
dieser Teil dreht sich mit der Hantel herum.

\begin{center}
\begin{tikzpicture}[scale=1.15]
  \draw[kldim,dashed] (0,-2.0)--(0,2.0);
  \draw[vec,klink] (0,1.75)--(0,2.45) node[right]{$\vec\omega$};
  \draw[kldim,-{Stealth[length=1.6mm]}] ([shift={(0,2.05)}]200:.38 and .12) arc (200:500:.38 and .12);
  \draw[klhi,line width=1.3pt] (-1.03,-1.47)--(1.03,1.47);
  \fill[klhi] (1.03,1.47) circle (4.5pt); \fill[klhi] (-1.03,-1.47) circle (4.5pt);
  \node[right=5pt] at (1.03,1.47){$m$}; \node[left=5pt] at (-1.03,-1.47){$m$};
  \draw[kldim] (0,.6) arc (90:55:.6); \node[kldim] at (.22,.85){$\theta$};
  \draw[vec,klhi,line width=1.1pt] (0,0)--(-1.38,.97) node[above left]{$\vec L$};
  \draw[kldim] (-.229,.161)--(-.068,.390)--(.161,.229);
  \node[kldim,anchor=west] at (1.6,-.4){\begin{minipage}{40mm}\fell\itshape\fontsize{10.5}{13}\selectfont $\vec L\perp$ Stab, nicht $\parallel\vec\omega$.\\ Dreht die Hantel, läuft $\vec L$ auf einem Kegel um die Achse.\end{minipage}};
\end{tikzpicture}
\end{center}
\plotcaption{Momentaufnahme. Die Spitze von $\vec L$ wandert, also ist $\dot{\vec L}\ne0$, obwohl $\omega$ konstant ist. Ein sich änderndes
$\vec L$ braucht ein Drehmoment, und das liefern die Lager der Achse. Das ist eine Unwucht; beim Reifenwechsel wird sie mit kleinen
Bleigewichten weggewuchtet.}

Intuition, im mitdrehenden System: Die Fliehkraft zieht beide Massen nach außen, die obere nach rechts, die untere nach links. Die
beiden Kräfte bilden ein Kräftepaar, das den Stab waagrecht stellen will. Die Lager halten dagegen. Die genaue Größe:
$|\vec D| = |\dot{\vec L}| = \omega|L_x| = 2m\omega^2a^2\sin\theta\cos\theta$.

\subsection{Der allgemeine Fall}
Dreht sich der Körper um einen festen Punkt (oder betrachtest du die Drehung um $S$), ist $\vec v = \vec\omega\times\vec r$ mit $\vec r$
von diesem Punkt aus. Mit $\vec A\times(\vec B\times\vec C) = \vec B(\vec A\cdot\vec C)-\vec C(\vec A\cdot\vec B)$:
\[
  \vec L = \int\vec r\times(\vec\omega\times\vec r)\,\dd m = \int\big[r^2\vec\omega - (\vec r\cdot\vec\omega)\,\vec r\big]\,\dd m,
  \qquad L_i = \sum_j\Big[\int(r^2\delta_{ij}-x_ix_j)\,\dd m\Big]\omega_j .
\]

\begin{definition}{Trägheitstensor}{inertia tensor}
Bezüglich eines Punktes (Schwerpunkt oder fester Drehpunkt) ist $I_{ij} = \int(r^2\delta_{ij}-x_ix_j)\,\dd m$, ausgeschrieben
\[
  I = \begin{pmatrix}
    \int(y^2+z^2)\,\dd m & -\int xy\,\dd m & -\int xz\,\dd m\\[2pt]
    -\int xy\,\dd m & \int(x^2+z^2)\,\dd m & -\int yz\,\dd m\\[2pt]
    -\int xz\,\dd m & -\int yz\,\dd m & \int(x^2+y^2)\,\dd m
  \end{pmatrix}.
\]
Dann gilt $\vec L = I\vec\omega$, also $L_i = \sum_jI_{ij}\,\omega_j$, und $E_{\mathrm{rot}} = \frac12\vec\omega\cdot\vec L = \frac12\,\vec\omega^{\mathrm T}I\,\vec\omega$.
\end{definition}

Lies die Matrix so. Auf der Diagonale stehen die Trägheitsmomente um die Koordinatenachsen: $y^2+z^2$ ist das Abstandsquadrat von der
$x$-Achse. Daneben stehen die Deviationsmomente \de{(products of inertia)}. Sie messen, wie schief die Masse zu den Achsen liegt, und
sie sind genau das, was $\vec L$ von $\vec\omega$ wegdreht. $I$ ist symmetrisch, $I_{ij} = I_{ji}$, weil $x_ix_j = x_jx_i$.

Die Energie in einer Zeile: $(\vec\omega\times\vec r)^2 = \omega^2r^2-(\vec\omega\cdot\vec r)^2 = \vec\omega\cdot\big[r^2\vec\omega-(\vec r\cdot\vec\omega)\vec r\big]$,
integriert $E_{\mathrm{rot}} = \frac12\vec\omega\cdot\vec L$. Ausgeschrieben, wie bei dir:
$\frac12\big(I_{11}\omega_1^2+I_{22}\omega_2^2+I_{33}\omega_3^2+2I_{12}\omega_1\omega_2+2I_{13}\omega_1\omega_3+2I_{23}\omega_2\omega_3\big)$.
Und für $\vec\omega = \omega\hat n$ ist $E_{\mathrm{rot}} = \frac12(\hat n^{\mathrm T}I\hat n)\,\omega^2$: Das Trägheitsmoment um \emph{jede}
Achse $\hat n$ durch den Punkt ist $I_n = \hat n^{\mathrm T}I\hat n = \sum_{ij}I_{ij}n_in_j$. Der Tensor kennt sie alle.

\begin{panel}[unbreakable]{Check}{deine Mitschrift}{}
In der Zeile für $L_1$ steht bei dir $+\,xy\,\omega_2 + xz\,\omega_3$. Aus $-(\vec r\cdot\vec\omega)\,\vec r$ kommt dort ein Minus:
\[
  L_1 = \rho\int\big[(y^2+z^2)\,\omega_1 - xy\,\omega_2 - xz\,\omega_3\big]\,\dd^3V .
\]
Deine Definition $I_{12} = -\int x_1x_2\,\dd m$ weiter unten ist richtig, und sie passt nur mit diesem Minus. Ein Vorzeichen, kein
Denkfehler; aber genau so eines kostet in der Klausur den ganzen Tensor.
\end{panel}

\subsection{Hauptachsen}

\begin{theorem}{Hauptachsen}{principal axes}
$I$ ist reell und symmetrisch. Nach dem Spektralsatz (Lineare Algebra) gibt es drei zueinander senkrechte Einheitsvektoren
$\hat e_1,\hat e_2,\hat e_3$ und Zahlen $I_1,I_2,I_3\ge0$ mit $I\hat e_k = I_k\hat e_k$. In diesen Achsen ist
\[
  I = \begin{pmatrix}I_1&0&0\\0&I_2&0\\0&0&I_3\end{pmatrix},\qquad \vec L = (I_1\omega_1,\ I_2\omega_2,\ I_3\omega_3),\qquad
  E_{\mathrm{rot}} = \tfrac12\big(I_1\omega_1^2+I_2\omega_2^2+I_3\omega_3^2\big).
\]
$\vec L\parallel\vec\omega$ gilt genau dann, wenn $\vec\omega$ ein Eigenvektor von $I$ ist.
\end{theorem}

Die $I_k$ heißen Hauptträgheitsmomente. Sie sind $\ge0$, weil $\hat n^{\mathrm T}I\hat n = \int r_\perp^2\,\dd m\ge0$. Bei der Hantel ist die
Achse kein Eigenvektor, darum die Unwucht. Eine Achse, um die sich ein Körper ohne Lagerkräfte drehen kann, heißt freie Achse; das sind
genau die Hauptachsen durch den Schwerpunkt.

\textbf{Hauptachsen ohne Rechnung, aus der Symmetrie.} Ist $z\mapsto-z$ eine Spiegelsymmetrie des Körpers, hat jedes $\dd m$ bei
$(x,y,z)$ einen Partner bei $(x,y,-z)$. In $I_{xz} = -\int xz\,\dd m$ und $I_{yz}$ heben sich die Partner auf: Die Normale einer
Spiegelebene ist eine Hauptachse. Ebenso jede Drehsymmetrieachse. Dein $I_{xy} = 0$ beim Zylinder, mit dem
$\int_0^{2\pi}\cos\varphi\sin\varphi\,\dd\varphi = 0$, ist kein Zufall: Es ist die Symmetrie, die in der Rechnung spricht.

Drei Sorten Körper. Alle $I_k$ gleich: Kugelkreisel (Kugel, Würfel). Zwei gleich: symmetrischer Kreisel (Zylinder, Kegel, näherungsweise die
Erde). Alle verschieden: unsymmetrischer Kreisel (Buch, Handy, Tennisschläger).

\why{Warum hat ein Würfel um jede Achse durch seinen Mittelpunkt dasselbe Trägheitsmoment, sogar um die Raumdiagonale?}{Die drei
Spiegelebenen machen alle Deviationsmomente null, und die Drehungen um $90^\circ$ vertauschen die Achsen, also $I_1 = I_2 = I_3 = \frac16Ma^2$.
Dann ist $I$ das $\frac16Ma^2$-fache der Einheitsmatrix, und $\hat n^{\mathrm T}I\hat n = \frac16Ma^2$ für jedes $\hat n$. Ein Würfel dreht
sich wie eine Kugel. Er sieht nur nicht so aus.}

Zwei Werkzeuge noch. \emph{Steiner für den Tensor:} Ist $\vec d$ der Vektor von $A$ nach $S$, dann
$I^{(A)}_{ij} = I^{(s)}_{ij} + M(d^2\delta_{ij}-d_id_j)$; die Diagonale davon ist der Steiner aus Teil IV. \emph{Die Dreiecksungleichung:}
$I_1+I_2 = I_3 + 2\int x_3^2\,\dd m\ge I_3$, mit Gleichheit genau für flache Körper. Das ist der Satz der senkrechten Achsen, von oben
gesehen. Trägst du für jede Richtung $\hat n$ den Punkt $\hat n/\sqrt{I_n}$ ab, entsteht die Fläche $\vec x^{\mathrm T}I\vec x = 1$, das
Trägheitsellipsoid, mit den Hauptachsen als Achsen. Beim Würfel ist es eine Kugel.

% ======================================================================
\section{Die zwei Gleichungen}

\opening{J}{etzt die} Dynamik, und sie ist kürzer, als du fürchtest. Schreib Newton für jedes Teilchen hin,
$m_i\ddot{\vec r}_i = \vec F_i^{\mathrm{ext}} + \sum_{j\ne i}\vec F_{ij}$, und summiere.

\emph{Über alle Teilchen:} Die inneren Kräfte heben sich paarweise auf ($\vec F_{ij} = -\vec F_{ji}$, actio gleich reactio). Es bleibt
$M\ddot{\vec R}_s = \vec F\ext$, der Schwerpunktsatz.

\emph{Mit $\vec r_i\times$ davor:} $\dot{\vec L} = \sum_i\vec r_i\times m_i\ddot{\vec r}_i$ (denn $\dot{\vec r}_i\times\dot{\vec r}_i = 0$).
Ein Paar innerer Kräfte gibt $\vec r_i\times\vec F_{ij}+\vec r_j\times\vec F_{ji} = (\vec r_i-\vec r_j)\times\vec F_{ij}$, und das ist null,
wenn die Kraft entlang der Verbindungslinie wirkt. Es bleibt $\dot{\vec L} = \vec D\ext$, der Drallsatz, um einen raumfesten Punkt.

\emph{Zerlegt in Schwerpunkt und Drehung um ihn:} Setz $\vec v_i = \vec v_s+\vec\omega\times\vec r_{si}$ ein. Bei der Energie
\[
  E = \sum_i\tfrac12m_i\big(\vec v_s+\vec\omega\times\vec r_{si}\big)^2
    = \tfrac12Mv_s^2 + \vec v_s\cdot\Big(\vec\omega\times\sum_im_i\vec r_{si}\Big) + \tfrac12\sum_im_i(\vec\omega\times\vec r_{si})^2 ,
\]
und der Kreuzterm stirbt am Arbeitspferd, zum dritten Mal. Beim Drehimpuls genauso.

\begin{keystone}{Die zwei Gleichungen}{equations of motion}
\[
  M\,\dot{\vec v}_s = \vec F\ext, \qquad \dot{\vec L} = \vec D\ext, \qquad \vec L = I\vec\omega .
\]
Die zweite gilt um einen raumfesten Punkt oder um den Schwerpunkt, auch wenn dieser beschleunigt. Dazu
\[
  E_{\mathrm{kin}} = \tfrac12Mv_s^2 + \tfrac12\,\vec\omega\cdot I_s\vec\omega, \qquad
  \vec L_O = M\,\vec R_s\times\vec v_s + I_s\vec\omega ,
\]
mit $I_s$ um den Schwerpunkt: Bahnanteil plus Eigendrehung. Sechs Gleichungen für sechs Freiheitsgrade.

\smallskip\noindent{\color{kldim}\pix{Beweis in einem Atemzug}}\quad Newton für jedes Teilchen, summiert; innere Kräfte heben sich
paarweise weg (actio gleich reactio, Zentralkräfte); bei $E$ und $\vec L$ fallen die Kreuzterme wegen $\sum m_i\vec r_{si} = 0$.
\end{keystone}

\why{Warum gilt der Drallsatz um den Schwerpunkt, auch wenn der Schwerpunkt beschleunigt, aber nicht um jeden anderen bewegten
Punkt?}{Rechnung: $\vec L_O = M\vec R_s\times\vec v_s+\vec L_s$ ableiten gibt $M\vec R_s\times\dot{\vec v}_s+\dot{\vec L}_s$, und
$\vec D_O = \sum(\vec R_s+\vec r_{si})\times\vec F_i = \vec R_s\times\vec F+\vec D_s$. Mit $M\dot{\vec v}_s = \vec F$ bleibt
$\dot{\vec L}_s = \vec D_s$. Bild: Im beschleunigten System eines Punktes $P$ wirkt die Scheinkraft $-M\vec a_P$, und zwar im Schwerpunkt.
Um $S$ hat sie keinen Hebelarm, um jeden anderen Punkt schon. Nimm also $S$ oder einen raumfesten Punkt, sonst nichts.}

\subsection{Was erhalten bleibt}
\emph{Drehimpuls:} Ist $\vec D\ext = 0$ um einen Punkt, bleibt $\vec L$ um diesen Punkt konstant, als Vektor. \emph{Energie:} Bleibt
erhalten, wenn nur konservative Kräfte Arbeit leisten. Zwangskräfte einer reibungsfreien festen Achse und die Haftreibung beim Rollen
leisten keine (Teil VII).

\why{Eine Eisläuferin zieht in der Pirouette die Arme an. Der Drehimpuls bleibt, die Energie nicht. Woher kommt sie?}{$L = I\omega$
bleibt, also $E = L^2/(2I)$, und das steigt, wenn $I$ sinkt. Die Differenz leistet sie selbst. Im mitdrehenden System sieht man es: Die
Arme müssen gegen die Fliehkraft $m\omega^2r$ nach innen gezogen werden, und $\int m\omega^2r\,\dd r$ ist genau $\Delta E$. Drehimpuls
kann von innen nicht nachgefüllt werden, Energie schon.}

\begin{panel}{Rezept}{jede Aufgabe zum starren Körper}{}
\begin{enumerate}
  \item \textbf{Skizze.} Jede Kraft mit ihrem Angriffspunkt. Beim starren Körper zählt der Ort, nicht nur der Betrag.
  \item \textbf{Zwangsbedingungen.} Feste Achse? Rollen, $v_s = \omega R$ und $\dot v_s = \dot\omega R$? Seil ohne Schlupf, $a = \dot\omega r$?
  \item \textbf{Welcher Satz.} Kräfte oder Beschleunigungen gesucht: die zwei Gleichungen. Nur Geschwindigkeiten zwischen zwei Lagen:
        Energie, falls nichts Energie verheizt. Stoß, Zusammenziehen, unbekannte Kraft in einem Gelenk: Drehimpulserhaltung um dieses Gelenk.
  \item \textbf{Bezugspunkt.} $S$ oder ein raumfester Punkt, am besten dort, wo Unbekannte angreifen.
  \item \textbf{Zählen.} So viele Gleichungen wie Unbekannte?
  \item \textbf{Rückschau.} Einheiten. $I\to0$ muss das Ergebnis für Massenpunkte geben. Grenzfälle.
\end{enumerate}
\end{panel}

\begin{example}{Rolle mit Masse}
Zwei Massen $m_1<m_2$ hängen an einem Seil über einer Rolle (Vollzylinder, Masse $M$, Radius $R$, reibungsfrei gelagert). Das Seil rutscht
nicht. Wie groß ist die Beschleunigung?
\step{Verstehen} Drei Körper, Unbekannte $a$, $T_1$, $T_2$. Die Seilkräfte links und rechts sind \emph{verschieden}: Wären sie gleich,
hätte die Rolle kein Drehmoment und könnte nicht schneller werden.
\step{Plan} Newton für jede Masse, $D = I\dot\omega$ für die Rolle, dazu die Zwangsbedingung $a = \dot\omega R$.
\step{Rechnung} $m_2a = m_2g-T_2$, \ $m_1a = T_1-m_1g$, \ $(T_2-T_1)R = \frac12MR^2\cdot\frac aR$, also $T_2-T_1 = \frac12Ma$. Die ersten
zwei addiert: $(m_1+m_2)a = (m_2-m_1)g-(T_2-T_1)$, also
\[
  a = \frac{(m_2-m_1)\,g}{m_1+m_2+\frac12M} .
\]
\step{Rückschau} $M\to0$ gibt die ideale Atwood-Maschine. Die Rolle wirkt wie eine Zusatzmasse $\frac12M$, die mitbeschleunigt werden muss,
an der aber keine Schwerkraft zieht (die trägt das Lager). Das $\frac12$ ist ihr $k = I/(MR^2)$.
\end{example}

\begin{shortexample}{Kugel bleibt in einem Stab stecken}
Ein Stab (Masse $M$, Länge $\ell$) hängt reibungsfrei an seinem oberen Ende. Eine Kugel (Masse $m$, Geschwindigkeit $v$, waagrecht) trifft
ihn am unteren Ende und bleibt stecken. Winkelgeschwindigkeit direkt danach?
\step{Plan} Impuls? Nein: Das Lager stößt während des Stoßes mit einer unbekannten Kraft zurück. Energie? Nein: Steckenbleiben verheizt
Energie. Aber um das Lager hat die Lagerkraft keinen Hebelarm, und die Schwerkraft ist während des kurzen Stoßes vernachlässigbar.
Also bleibt der Drehimpuls um das Lager erhalten.
\step{Rechnung} Vorher $L = mv\ell$. Nachher $L = \big(\frac13M\ell^2+m\ell^2\big)\omega$. Also
\[
  \omega = \frac{3mv}{(M+3m)\,\ell} .
\]
\step{Rückschau} $M\to0$: $\omega = v/\ell$, die Kugel läuft einfach weiter auf dem Kreis. Einheiten: $1/\mathrm s$.
\end{shortexample}

% ======================================================================
\section{Rollen}

\opening{E}{in Rad} rollt ohne zu rutschen, wenn der Auflagepunkt $P$ dieselbe Geschwindigkeit hat wie die Straße: null. Mit dem Satz
aus Teil II und $\vec r_{sP}$ vom Mittelpunkt nach unten:
\[
  \vec v_P = \vec v_s + \vec\omega\times\vec r_{sP} = 0 \quad\Longrightarrow\quad v_s = \omega R .
\]
Das ist die Rollbedingung \de{(rolling condition)}. Oben zeigt $\vec\omega\times\vec r$ nach vorn statt nach hinten, also
$v_{\mathrm{oben}} = v_s+\omega R = 2v_s$. Und weil $P$ gerade ruht, darfst du (Warum-Frage in Teil II) die ganze Bewegung als reine
Drehung um $P$ lesen: $\vec v_i = \vec\omega\times(\vec r_i-\vec r_P)$. Jeder Punkt bewegt sich senkrecht zu seiner Verbindung mit $P$,
mit $\omega$ mal Abstand zu $P$.

\begin{center}
\begin{tikzpicture}[scale=1.5]
  \fill[pattern={dither25},pattern color=kllow] (-1.9,-1.2) rectangle (2.7,-1); \draw[kldim] (-1.9,-1)--(2.7,-1);
  \draw[klink,line width=.9pt] (0,0) circle (1);
  \foreach \p in {(1,0),(-1,0)} \draw[kldim,dashed,line width=.4pt] (0,-1)--\p;
  \draw[kldim,dashed,line width=.4pt] (0,-1)--(1.2,1);
  \draw[vec,klhi,line width=1pt] (0,1)--(1.2,1) node[right]{$2v_s$};
  \draw[vec,klhi,line width=1pt] (0,0)--(.6,0) node[above=2pt,lab]{$v_s$};
  \draw[vec,klhi,line width=1pt] (1,0)--(1.6,-.6) node[right]{$\sqrt2\,v_s$};
  \draw[vec,klhi,line width=1pt] (-1,0)--(-.4,.6) node[above=1pt,lab]{$\sqrt2\,v_s$};
  \fill[klhi] (0,-1) circle (2pt) node[left=8pt,lab]{$P$: \ $v = 0$};
  \fill[klhi] (0,0) circle (1.5pt) (0,1) circle (1.5pt) (1,0) circle (1.5pt) (-1,0) circle (1.5pt);
  \draw[vec,klink] (225:.55) arc (225:135:.55); \node[klink] at (-.27,0){$\omega$};
\end{tikzpicture}
\end{center}
\plotcaption{Das Geschwindigkeitsprofil entlang des senkrechten Durchmessers ist eine Gerade durch $P$: unten $0$, Mitte $v_s$, oben $2v_s$.
Die seitlichen Punkte bewegen sich schräg, senkrecht zu ihrer Linie nach $P$. Das Rad dreht sich in jedem Moment um seinen Auflagepunkt.}

\anim{rollen}{Erst schieben, dann drehen, dann Pfeil an Pfeil: oben $2v_s$, unten $0$. Achte darauf, dass jeder Ergebnispfeil
senkrecht auf seiner Linie zum Auflagepunkt steht. Am Ende rollt das Rad, und ein Randpunkt zeichnet eine Zykloide mit einer Spitze dort,
wo er kurz stillsteht.}

Die Energie zweimal, als Probe. Als Schwerpunkt plus Drehung: $E = \frac12Mv_s^2+\frac12I_s\omega^2 = \frac12\big(M+I_s/R^2\big)v_s^2$. Als
reine Drehung um $P$, mit Steiner $I_P = I_s+MR^2$: $E = \frac12I_P\omega^2 = \frac12(I_s+MR^2)\,v_s^2/R^2$. Dasselbe.

\subsection{Die schiefe Ebene}
Schreib $I_s = kMR^2$; die reine Zahl $k$ sagt nur, wie die Masse verteilt ist: Ring $1$, Hohlkugel $\frac23$, Vollzylinder $\frac12$,
Vollkugel $\frac25$. Rollt der Körper aus der Ruhe die Höhe $h$ hinunter, gibt der Energiesatz
\[
  Mgh = \tfrac12(1+k)\,Mv_s^2 \quad\Longrightarrow\quad v_s = \sqrt{\frac{2gh}{1+k}},\qquad a = \frac{g\sin\alpha}{1+k}.
\]
Probe mit Kräften. Haftreibung $f$ hangaufwärts: $Ma = Mg\sin\alpha-f$, und um $S$: $fR = I_s\dot\omega = kMR^2\cdot a/R$, also $f = kMa$.
Eingesetzt: $a(1+k) = g\sin\alpha$. Gleich. Nebenbei: Rollen ohne Rutschen braucht $f\le\mu_0Mg\cos\alpha$, also $\tan\alpha\le\mu_0(1+k)/k$.

\begin{center}
\begin{tikzpicture}
\begin{axis}[klplot,xbar,width=.78\linewidth,height=48mm,xmin=0,xmax=1.12,ymin=.4,ymax=5.6,
  ytick={1,2,3,4,5},yticklabels={Ring,Hohlkugel,Vollzylinder,Vollkugel,{Klotz, gleitet}},
  yticklabel style={font=\fell\itshape\fontsize{10}{11}\selectfont,text=klink},
  xtick={0,.5,1},xticklabels={0,$\frac12$,1},
  xlabel={$a\,/\,g\sin\alpha = 1/(1+k)$},bar width=8pt,axis lines*=left]
  \addplot[draw=klhi,line width=.6pt,pattern={dither50},pattern color=klhi] coordinates {(.5,1) (.6,2) (.6667,3) (.7143,4) (1,5)};
  \node[anchor=west,text=klhi] at (axis cs:.5,1){\ $\frac12$};
  \node[anchor=west,text=klhi] at (axis cs:.6,2){\ $\frac35$};
  \node[anchor=west,text=klhi] at (axis cs:.6667,3){\ $\frac23$};
  \node[anchor=west,text=klhi] at (axis cs:.7143,4){\ $\frac57$};
  \node[anchor=west,text=klhi] at (axis cs:1,5){\ $1$};
\end{axis}
\end{tikzpicture}
\end{center}
\plotcaption{Das Rennen: Vollkugel vor Vollzylinder vor Hohlkugel vor Ring, bei jeder Masse und jedem Radius. Nur ein reibungsfrei gleitender
Klotz ist schneller, denn er steckt nichts in die Drehung.}

\why{Warum entscheiden weder Masse noch Radius das Rennen?}{$a = g\sin\alpha/(1+k)$ hängt nur von $k = I_s/(MR^2)$ ab, und $k$ ist eine
reine Zahl, die nur sagt, \emph{wie} die Masse verteilt ist: innen oder außen. Jedes Kilogramm bekommt dieselbe Energie $gh$. Der Ring steckt
davon die Hälfte in die Drehung ($k/(1+k) = \frac12$), die Vollkugel nur $\frac27$. Was in der Drehung steckt, fehlt der Geschwindigkeit.}

\why{Warum leistet die Haftreibung beim Rollen keine Arbeit, obwohl sie eine Kraft ist?}{Leistung ist $P = \vec F\cdot\vec v$, mit der
Geschwindigkeit des \emph{Angriffspunkts}. Die Reibung greift am Auflagepunkt an, und der ruht: $P = 0$. Sie schiebt nur Energie um, von
der Translation in die Drehung. Darum darfst du beim Rollen ohne Gleiten den Energiesatz benutzen. Gleitreibung dagegen greift an einem
rutschenden Punkt an und verheizt.}

\begin{example}{Vom Gleiten zum Rollen}
Eine Billardkugel (Masse $M$, Radius $R$) wird so angestoßen, dass sie mit $v_0$ ohne Drehung über den Tisch gleitet. Die Gleitreibung bremst
sie und dreht sie an, bis sie rollt. Wie schnell rollt sie dann?
\step{Verstehen} Die Reibung greift am Auflagepunkt an, ist unbekannt (kein $\mu$ gegeben) und verheizt Energie. Gesucht: $v$ am Ende.
\step{Plan} Energie geht nicht. Kraft und Moment gingen, mit $\mu$, das sich herauskürzen müsste. Kürzer: Drehimpuls um einen raumfesten
Punkt $Q$ auf der Tischplatte. Die Reibung wirkt entlang der Tischplatte, durch $Q$ hindurch gedacht, also ohne Hebelarm. Gewicht und
Normalkraft sind gleich groß, entgegengesetzt und liegen auf derselben Senkrechten; ihre Momente heben sich auf. Also: $\vec L_Q$ bleibt.
\step{Rechnung} Mit der Keystone-Zerlegung ist $L_Q = Mv_sR+I_s\omega$ (Bahnanteil plus Eigendrehung; beide drehen im selben Sinn).
Anfangs $Mv_0R$. Am Ende rollt sie, $\omega = v/R$: $MvR+\frac25MR^2\cdot\frac vR = \frac75MvR$. Also
\[
  v = \tfrac57\,v_0 .
\]
\step{Rückschau} Das $\mu$ bestimmt nur, \emph{wie lange} sie rutscht, nicht wie schnell sie am Ende rollt. Allgemein $v = v_0/(1+k)$. Energie:
am Ende $\frac12\cdot\frac75Mv^2 = \frac57\cdot\frac12Mv_0^2$. Zwei Siebtel sind Wärme geworden.
\end{example}

% ======================================================================
\section{Das physikalische Pendel}

\opening{H}{äng einen} beliebigen Körper an einer waagrechten Achse $A$ auf, mit dem Schwerpunkt im Abstand $d$ darunter, und lenk ihn um
$\varphi$ aus. Die Schwerkraft greift in $S$ an und hat den Hebelarm $d\sin\varphi$: $D_z = -Mgd\sin\varphi$. Feste Achse, also
\[
  I_A\,\ddot\varphi = -Mgd\sin\varphi \quad\stackrel{\varphi\ \mathrm{klein}}{\Longrightarrow}\quad
  \ddot\varphi = -\frac{Mgd}{I_A}\,\varphi,\qquad T = 2\pi\sqrt{\frac{I_A}{Mgd}},\qquad I_A = I_s+Md^2 .
\]
Ein mathematisches Pendel der Länge $\ell_r = I_A/(Md)$, die reduzierte Pendellänge, hat dieselbe Periode. Der Stab am Ende aufgehängt:
$I_A = \frac13M\ell^2$, $d = \frac\ell2$, also $\ell_r = \frac23\ell$ und $T = 2\pi\sqrt{2\ell/(3g)}$. Kürzer als $\ell$, weil ein Teil der
Masse nah an der Achse hängt und schneller schwingen will.

\begin{center}
\begin{tikzpicture}
\begin{axis}[klplot,width=.84\linewidth,height=52mm,xmin=0,xmax=4.2,ymin=0,ymax=2.2,xlabel={$d\,/\,k_s$},ylabel={$T/T_{\min}$},
  ylabel style={rotate=-90},xtick={0,1,2,3,4},ytick={0,1,2},domain=.12:4.2,samples=160]
  \addplot[klhi]{sqrt((1/x+x)/2)} node[pos=.03,right=3pt]{$T(d)$};
  \addplot[kldim,dashed,domain=0:4.2]{sqrt(x/2)};
  \draw[kldim,dashed,line width=.5pt] (axis cs:1,0)--(axis cs:1,1);
  \node[anchor=north west,text=kldim] at (axis cs:2.3,1.0){\fell\itshape math.\ Pendel, $\ell = d$};
  \node[anchor=west,text=klink] at (axis cs:1,.45){$d = k_s$};
\end{axis}
\end{tikzpicture}
\end{center}
\plotcaption{Periode gegen Aufhängeabstand, mit dem Trägheitsradius $k_s = \sqrt{I_s/M}$. Bei $d\to0$ wird $T$ unendlich (im Schwerpunkt
aufgehängt: indifferent, Teil III). Weit draußen wird es ein mathematisches Pendel. Dazwischen ein Minimum bei $d = k_s$.}

Woher das Minimum: $T^2\propto I_A/(Md) = k_s^2/d + d$, und $k_s^2/d+d$ ist am kleinsten bei $d = k_s$. Jede Periode über dem Minimum wird
zweimal getroffen, bei $d$ und bei $d' = k_s^2/d$, und dann ist $d+d' = \ell_r$.

\aside{Kater hat 1817 genau damit $g$ gemessen: zwei Schneiden mit gleicher Periode, ihr Abstand ist $\ell_r$. Kein $I$ nötig. Ich war beeindruckt; das passiert selten.}

% ======================================================================
\section{Kreisel}

\opening{H}{ier zahlt} sich der wichtigste Satz des Kapitels aus: Ein Drehmoment ändert nicht die Lage, es ändert $\vec L$, und zwar in
\emph{seine} Richtung, $\dd\vec L = \vec D\,\dd t$.

\subsection{Präzession}
Ein Rad dreht sich schnell um seine Achse, Drehimpuls $\vec L = I_3\omega_3\,\hat e$ entlang der Achse. Die Achse ist waagrecht und
nur an einem Ende aufgehängt, der Schwerpunkt im Abstand $d$ davon. Die Schwerkraft gibt um den Aufhängepunkt $\vec D = \vec r\times M\vec g$:
waagrecht und senkrecht zu $\vec L$. In der Zeit $\dd t$ kommt $\dd\vec L = \vec D\,\dd t$ dazu, seitwärts. $\vec L$ wird nicht länger,
nicht kürzer, und kippt nicht nach unten: Es dreht sich in der Waagrechten. Die Achse folgt $\vec L$.

\begin{center}
\begin{tikzpicture}[scale=1.3]
  \draw[kldim,line width=.5pt] (-.6,2.2)--(.6,2.2);
  \foreach \x in {-.5,-.25,...,.5} \draw[kllow] (\x,2.2)--++(.15,.15);
  \draw[klink,line width=.6pt] (0,2.2)--(0,0);
  \draw[vec,klink] (-.45,.5)--(-.45,1.5) node[left]{$\vec\Omega$};
  \draw[klink,line width=1.4pt] (0,0)--(2.75,0);
  \fill[klhi] (0,0) circle (1.8pt);
  \draw[klhi,line width=.9pt,fill=klpage,fill opacity=.85] plot[smooth cycle,samples=48,domain=0:360,variable=\t]
    ({2.2+.5*.9*cos(\t)},{.35*.9*cos(\t)+.9*sin(\t)});
  \draw[klink,line width=1.4pt] (2.2,0)--(2.75,0);
  \fill[klhi] (2.2,0) circle (1.8pt);
  \draw[vec,klhi,line width=1.1pt] (2.75,0)--(3.95,0) node[above left=1pt]{$\vec L$};
  \draw[vec,klhi] (2.2,0)--(2.2,-1.35) node[below]{$M\vec g$};
  \draw[vec,klink,line width=1pt] (3.95,0)--(4.55,.42) node[right]{$\dd\vec L = \vec D\,\dd t$};
  \draw[kldim,dashed,line width=.5pt] plot[smooth,samples=20,domain=-14:24,variable=\t] ({3.95*cos(\t)+.5*3.95*sin(\t)},{.35*3.95*sin(\t)});
  \draw[kldim,line width=.4pt] (0,-.12)--(0,-.3) (2.2,-.12)--(2.2,-.3); \draw[kldim,<->,line width=.4pt] (0,-.22)--(2.2,-.22);
  \node[kldim,below] at (1.1,-.22){$d$};
\end{tikzpicture}
\end{center}
\plotcaption{Das Drehmoment der Schwerkraft zeigt waagrecht in die Tiefe des Bildes, senkrecht zu $\vec L$. Also wandert die Spitze von $\vec L$
in die Tiefe, auf dem gestrichelten Kreis. Von oben gesehen dreht sich die Achse gegen den Uhrzeigersinn, mit $\vec\Omega$ nach oben.}

In der Zeit $\dd t$ dreht sich $\vec L$ um den Winkel $\dd\phi = |\dd\vec L|/L = D\,\dd t/L$. Die Präzessionsgeschwindigkeit:
\[
  \Omega = \frac DL = \frac{Mgd}{I_3\,\omega_3}, \qquad\text{vektoriell}\quad \vec\Omega\times\vec L = \vec D .
\]
Schneller drehen heißt langsamer präzedieren. Die Rechnung nimmt an, dass $\vec L$ fast ganz aus der Eigendrehung besteht, also
$\Omega\ll\omega_3$ (schneller Kreisel); der kleine Drehimpuls der Präzessionsbewegung selbst ist vernachlässigt.

Und das Loslassen? Die Schwerkraft hat um den Aufhängepunkt kein Drehmoment um die Senkrechte ($\vec D\perp\vec g$), also bleibt
$L_z$ erhalten. Am Anfang ist $L_z = 0$. Damit die Präzession (die ein wenig Drehimpuls nach oben trägt) anlaufen kann, muss die Achse
ein Stück nach unten nicken, sodass die Eigendrehung ein wenig nach unten zeigt und das ausgleicht. Sie nickt über das Ziel hinaus und
zurück: die Nutation \de{(nutation)}. Reibung dämpft das Nicken weg, die ruhige Präzession bleibt.

\subsection{Die Euler-Gleichungen}
Für die volle Dynamik rechnet man im mitdrehenden System der Hauptachsen, wo $I$ konstant und diagonal ist. Für jeden Vektor gilt
\[
  \Big(\frac{\dd\vec A}{\dd t}\Big)_{\text{Raum}} = \Big(\frac{\dd\vec A}{\dd t}\Big)_{\text{Körper}}+\vec\omega\times\vec A
\]
(für einen im Körper festen Vektor ist das $\vec v = \vec\omega\times\vec r$). Mit $\vec A = \vec L = (I_1\omega_1,I_2\omega_2,I_3\omega_3)$ wird
$\dot{\vec L} = \vec D$ zu den Euler-Gleichungen
\begin{align*}
  I_1\dot\omega_1-(I_2-I_3)\,\omega_2\omega_3 &= D_1,\\
  I_2\dot\omega_2-(I_3-I_1)\,\omega_3\omega_1 &= D_2,\\
  I_3\dot\omega_3-(I_1-I_2)\,\omega_1\omega_2 &= D_3 .
\end{align*}

\begin{theorem}{Stabile und instabile Drehachsen}{tennis racket theorem}
Ein kräftefreier Körper mit $I_1>I_2>I_3$ dreht sich stabil um die Achse des größten und um die des kleinsten Hauptträgheitsmoments,
instabil um die mittlere.
\end{theorem}
\begin{proof}[Beweis, linearisiert]
Drehung fast um Achse 2: $\omega_2\approx\Omega$, $\omega_1,\omega_3$ klein, $\vec D = 0$. Die erste und dritte Euler-Gleichung, Produkte
kleiner Größen weggelassen: $I_1\dot\omega_1 = (I_2-I_3)\,\Omega\,\omega_3$ und $I_3\dot\omega_3 = (I_1-I_2)\,\Omega\,\omega_1$. Ableiten und
einsetzen:
\[
  \ddot\omega_1 = \frac{(I_2-I_3)(I_1-I_2)}{I_1I_3}\,\Omega^2\,\omega_1 .
\]
Beide Klammern sind positiv: $\omega_1$ wächst exponentiell, die Drehung kippt. Für Achse 1 oder 3 ergibt dieselbe Rechnung einen
negativen Faktor, also eine Schwingung um die Achse: stabil.
\end{proof}

\aside{Dschanibekow sah es 1985 auf Saljut~7 an einer Flügelmutter, die sich in der Schwerelosigkeit alle paar Sekunden umdrehte. Die Gleichung stand seit 220 Jahren bei Euler.}Probier es mit einem Buch (mit Gummiband zu) oder deinem Handy: hochwerfen, um jede der drei Achsen drehen. Um die längste und die kürzeste
Achse dreht es ruhig. Um die mittlere überschlägt es sich jedes Mal mit einer halben Drehung. Beim symmetrischen Kreisel ($I_1 = I_2$)
gibt es keine mittlere Achse; wirfst du ihn schief, läuft seine Figurenachse auf einem Kegel um das feste $\vec L$, die kräftefreie Nutation.


% ======================================================================
\begin{triggers}
\trigger{Ausgedehnter Körper im homogenen Schwerefeld}{$M\vec g$ im Schwerpunkt angreifen lassen}
\trigger{Gleichgewicht, Kräfte an mehreren Stellen}{$\sum\vec F = 0$, $\sum\vec D = 0$ um den Punkt, an dem die meisten Unbekannten angreifen}
\trigger{Unbekannte Kraft in Gelenk, Lager oder Auflagepunkt}{Moment bzw. Drehimpuls um genau diesen Punkt}
\trigger{Drehung um eine feste Achse}{$D_z = I\dot\omega$ mit $I$ um genau diese Achse}
\trigger{Achse parallel zu einer Schwerpunktsachse}{Steiner, $I = I_s+Md^2$, immer über $S$}
\trigger{Integral für $I$}{Symmetrie zuerst: Zylinderkoordinaten, $I_x+I_y+I_z = 2\int r^2\dd m$, Spiegelebenen töten Deviationsmomente}
\end{triggers}

\begin{triggers}
\trigger{Rollen ohne Gleiten}{$v_s = \omega R$; Haftreibung leistet keine Arbeit, also Energiesatz mit $\frac12(1+k)Mv^2$}
\trigger{Seil über einer Rolle mit Masse}{verschiedene Seilkräfte, $(T_2-T_1)R = I\dot\omega$, $a = \dot\omega R$}
\trigger{Stoß gegen einen gelagerten Körper}{Drehimpuls um das Lager; Impuls und Energie meist nicht}
\trigger{Etwas zieht sich zusammen, kein äußeres Moment}{$I\omega$ konstant, $E = L^2/2I$ ändert sich}
\trigger{Gleiten geht in Rollen über}{Drehimpuls um einen Punkt auf der Unterlage}
\trigger{Lager werden belastet, Körper eiert}{$\vec L\nparallel\vec\omega$: keine Hauptachse}
\trigger{Drehmoment senkrecht zu einem großen $\vec L$}{Präzession, $\Omega = D/L$}
\trigger{Ausgedehnter Körper schwingt}{$T = 2\pi\sqrt{I_A/(Mgd)}$}
\end{triggers}

\needspace{.42\textheight}
\marginkael{proud}
\begin{deck}
\card{Def}{Trägheitstensor: Komponenten, und wie er $\vec L$ und $E$ gibt.}{$I_{ij} = \int(r^2\delta_{ij}-x_ix_j)\,\dd m$, \ $\vec L = I\vec\omega$, \ $E = \frac12\vec\omega\cdot I\vec\omega$}
\card{Def}{Steiner: $I_A = $ \{\{$I_s+Md^2$\}\}, wenn die Bezugsachse \{\{durch $S$\}\} geht und \{\{parallel\}\} zu $A$ ist.}{Von $A$ nach $B$ immer über $S$.}
\card{Warum}{Warum braucht ein Rad mit Unwucht Lagerkräfte, obwohl $\omega$ konstant ist?}{$\vec L\nparallel\vec\omega$ läuft mit dem Rad um. $\dot{\vec L}\ne0$ braucht ein Drehmoment: die Lager.}
\card{Warum}{Warum gewinnt die Vollkugel gegen den Ring, bei jeder Masse und jedem Radius?}{$a = g\sin\alpha/(1+k)$; $k = I_s/MR^2$ ist eine reine Formzahl: Kugel $\frac25$, Ring $1$.}
\card{Trigger}{Unbekannte Kraft in einem Gelenk oder Auflagepunkt. Wozu greifst du?}{Moment oder Drehimpuls um genau diesen Punkt.}
\end{deck}

\end{document}
